\section{Asymmetric Vaughan normal forms}
\label{sec:asymmetric-vaughan}

For integers $U,V\ge1$, Vaughan's identity in the sign convention used
here is
\begin{align}
 \sum_t\Lambda(t)F(t)
 ={}&\sum_{t\le U}\Lambda(t)F(t)
 -\sum_{\substack{\ell\le U\\d\le V}}
 \Lambda(\ell)\mu(d)\sum_{j\ge1}F(\ell dj)
 \notag\\
 &+\sum_{d\le V}\mu(d)
 \sum_{j\ge1}(\log j)F(dj)
 -\sum_{\substack{\ell>U\\k>1}}
 \Lambda(\ell)\beta_V(k)F(\ell k),
 \label{eq:general-Vaughan-identity}
\end{align}
where $\beta_V$ is defined in \eqref{eq:betaV-intro}.  Since
$\beta_V(k)=0$ for $1<k\le V$, the last term is supported on $k>V$.
The identity is classical \citep{Vaughan1977,HeathBrown1982}.

The TPC-15 proof did not require $U=V$.  Keeping the parameters free gives
the following useful invariance.

\begin{proposition}[Vaughan parameter invariance below the cutoff]
\label{prop:Vaughan-parameter-invariance}
Let $1\le U,V$ satisfy $UV\le R$.  Define
\begin{equation}
 \cH_{h,R;U,V}(X)=
 \sum_{\substack{\ell>U\\k>V}}
 \Lambda(\ell)\beta_V(k)r_R(\ell k+h)W(\ell k/X).
 \label{eq:general-hard-packet}
\end{equation}
Then for every $A>0$,
\begin{equation}
 \cP_h(X;W)=\mathfrak S(h)XI_0(W)
 -\cH_{h,R;U,V}(X)
 +O_{A,h,W,\delta}\bigl(X(\log X)^{-A}\bigr).
 \label{eq:general-parameter-reduction}
\end{equation}
Consequently any two admissible parameter pairs give hard packets that
differ by $O_A(X(\log X)^{-A})$.
\end{proposition}

\begin{proof}
TPC-15 proves the model cross-correlation
\begin{equation}
 \sum_t\Lambda(t)\Lambda_R(t+h)W(t/X)
 =\mathfrak S(h)XI_0(W)
 +O_A\bigl(X(\log X)^{-A}\bigr).
 \label{eq:imported-model-cross}
\end{equation}
Apply \eqref{eq:general-Vaughan-identity} to
$F(t)=r_R(t+h)W(t/X)$.  The support of $W$ makes the first packet zero
because $U\le R=o(X)$.  In the second packet the row modulus is
$m=\ell d\le UV\le R$, and its total coefficient at a fixed $m$ is at
most $\sum_{\ell\mid m}\Lambda(\ell)=\log m$.  The third packet has
$d\le V\le R$ and the standard logarithmic row multiplier.  The full
bounded Type-I row theorem of TPC-15 controls both packets with arbitrary
logarithmic saving.  The final term is
$-\cH_{h,R;U,V}$.  Combining with \eqref{eq:imported-model-cross} proves
the proposition.
\end{proof}

At the edge of \cref{prop:Vaughan-parameter-invariance}, the choice
$U=R,V=1$ gives $\beta_1(k)=1$.  The recentered theorem permits $U$ to
move much farther.

\begin{theorem}[Square-root Vaughan normal form]
\label{thm:square-root-normal-form}
Let $M$ be given by \eqref{eq:near-square-root-M} and define
\begin{equation}
 \cH^\sharp_{h,R;M}(X)=
 \sum_{\substack{\ell>M\\k>1}}
 \Lambda(\ell)r_R(\ell k+h)W(\ell k/X).
 \label{eq:Hsharp-definition}
\end{equation}
Then, for every $A>0$,
\begin{equation}
 \cP_h(X;W)=\mathfrak S(h)XI_0(W)
 -\cH^\sharp_{h,R;M}(X)
 +O_{A,h,W,\delta}\bigl(X(\log X)^{-A}\bigr).
 \label{eq:square-root-normal-form}
\end{equation}
In particular, when $h$ is an admissible even shift and $I_0(W)\ne0$,
\begin{equation}
 \cP_h(X;W)\sim\mathfrak S(h)XI_0(W)
 \quad\Longleftrightarrow\quad
 \cH^\sharp_{h,R;M}(X)=o(X).
 \label{eq:Hsharp-gate}
\end{equation}
\end{theorem}

\begin{proof}
Again split the target as $\Lambda(t+h)=\Lambda_R(t+h)+r_R(t+h)$ and
use \eqref{eq:imported-model-cross}.  Apply
\eqref{eq:general-Vaughan-identity} to the residual test function with
$U=M,V=1$.  The first packet vanishes by support.  The second packet is
\[
 -\sum_{\ell\le M}\Lambda(\ell)
 \sum_jr_R(\ell j+h)W(\ell j/X),
\]
which is $O_A(X(\log X)^{-A})$ by
\cref{cor:bounded-log-rows}.  The third is precisely the modulus-one
logarithmic row in \cref{lem:modulus-one-log-row}.  Finally,
$\beta_1(k)=1$, so the last packet is
$-\cH^\sharp_{h,R;M}$.  This proves \eqref{eq:square-root-normal-form}.
The equivalence follows because the displayed remainder is $o(X)$ and
the stated main coefficient is nonzero.
\end{proof}

\begin{corollary}[Equivalence with the TPC-15 packet]
\label{cor:two-hard-packets-equivalent}
With $U=V=\lfloor R^{1/2}\rfloor$,
\begin{equation}
 \cH^{\rm HB}_{h,R;U,V}(X)
 =\cH^\sharp_{h,R;M}(X)
 +O_A\bigl(X(\log X)^{-A}\bigr).
 \label{eq:packet-equivalence}
\end{equation}
\end{corollary}

This equality is an equality of total signed packets.  It does not match
individual factor pairs, and it must not be read as a positivity-preserving
transport from one dyadic scale to another.
